\documentclass{article}

\usepackage[preprint]{neurips_2026}

\usepackage[utf8]{inputenc}
\usepackage[T1]{fontenc}
\usepackage{hyperref}
\usepackage{url}
\usepackage{booktabs}
\usepackage{amsmath}
\usepackage{amssymb}
\usepackage{graphicx}
\usepackage{xcolor}

\hypersetup{colorlinks=true, linkcolor=black, citecolor=black, urlcolor=blue!60!black}

% Numbers generated from run artifacts; never typed by hand.
\newcommand{\OracleCeiling}{0.594}
\newcommand{\OracleSuccess}{0.500}
\newcommand{\BCProgress}{0.445}
\newcommand{\BCSuccess}{0.281}
\newcommand{\PPOProgress}{0.413}
\newcommand{\PPOSuccess}{0.250}
\newcommand{\RandomProgress}{0.146}
\newcommand{\BenchmarkTasks}{32}
\newcommand{\BCHeldOutAccuracy}{0.983}
\newcommand{\BCEpochs}{40}
\newcommand{\PPOIterations}{25}
\newcommand{\BCvsPPODelta}{+0.031}
\newcommand{\BCvsPPOCI}{[-0.008, +0.091]}
\newcommand{\BCvsPPOWLT}{5/1/26}
\newcommand{\CodecSteps}{16,452}
\newcommand{\CodecUnusable}{0}
\newcommand{\CodecWorstDrift}{0.0005}


\title{KAIROS: Measuring What Caps a Reinforcement-Learning Agent\\
       That Drives Parametric CAD}

\author{%
  Sunkalp Chandra \\
  \texttt{sunkalp.chandra@gmail.com} \\
}

\begin{document}
\maketitle

\begin{abstract}
We present KAIROS, an agent that builds parametric CAD models by emitting
structured operations into a live FreeCAD kernel from a natural-language
requirement, together with a benchmark that measures how far such an agent
actually gets. Our central finding is methodological: on a task of this shape,
\emph{the headline score is usually measuring the harness rather than the
policy}. We make that concrete. An oracle that replays the ground-truth expert
through the agent's own action codec scores \OracleCeiling{} progress and
\OracleSuccess{} full-build success --- not $1.0$ --- so every learned number
must be read against that ceiling rather than against perfection. Behavioural
cloning reaches \BCHeldOutAccuracy{} held-out next-action accuracy yet only
\BCProgress{} progress in closed loop, and a PPO stage anchored to it reaches
\PPOProgress{}; a paired bootstrap over identical tasks puts their difference at
\BCvsPPODelta{} with a 95\% interval of $\BCvsPPOCI$
(\BCvsPPOWLT{} win/loss/tie), so the benchmark \emph{cannot separate them}.
Ablations show the policy does read its requirement, and that much of its low
invalid-action rate belongs to the action mask rather than to the policy.
Along the way we report four defects that each produced a plausible, publishable,
wrong number, and the audits that now prevent them.
\end{abstract}

\section{Introduction}

Computer-aided design is an attractive target for sequential decision making. A
part is built by an ordered sequence of discrete, parameterised operations ---
sketch a profile, pad it, pocket a hole, fillet an edge --- and the result is
checkable: mass, wall thickness, hole count and angles are all measurable
against a written requirement. That combination, a discrete action space with a
programmatic reward, is rare outside of games.

It is also a setting where it is unusually easy to report a number that is not
about the policy. A CAD agent sits behind several layers of machinery --- an
action codec, a geometry kernel, a constraint checker, a requirement parser ---
and each layer can fail \emph{silently} and \emph{plausibly}. The agent still
runs, the kernel still returns valid solids, the checker still reports
satisfaction, and the resulting score looks like a result.

This paper is organised around that hazard. We describe the system
(Section~\ref{sec:system}), the benchmark (Section~\ref{sec:benchmark}), and the
results (Section~\ref{sec:results}). We then devote Section~\ref{sec:pitfalls}
to four defects we found in our own pipeline, each of which had already produced
a number we believed. We think this is the most transferable part of the work.

\paragraph{Contributions.}
\begin{enumerate}\itemsep2pt
\item A CAD environment exposing \textbf{38 structured operations} against a
      live FreeCAD kernel, with a JSON bridge that lets a PyTorch policy drive
      an interpreter that cannot import PyTorch.
\item \textbf{KAIROS-CAD-v1}, a benchmark of BUILD and COMPLETE($k$) tasks over
      a frozen three-way split, scored by a prefix-ordered milestone ladder
      rather than by success alone.
\item An \textbf{oracle-replay baseline} that measures the ceiling the action
      codec imposes. We argue this belongs in any agent benchmark whose action
      space is a lossy encoding of the demonstrations.
\item A catalogue of \textbf{silent-failure modes} in evaluating such agents,
      each with the measurement that exposes it, all now automated.
\end{enumerate}

\section{Related work}\label{sec:related}

\paragraph{Learning CAD construction.} DeepCAD~\citep{wu2021deepcad} and
Computer-Aided Design as Language~\citep{ganin2021computer} model construction
sequences generatively, and the Fusion 360 Gallery~\citep{willis2021fusion}
supplies human design sequences with a reconstruction environment. These
generate or reconstruct sequences; KAIROS instead executes actions against a
live kernel and scores the resulting solid against a written requirement, which
is what makes the harness itself a measurement problem.

\paragraph{Imitation and compounding error.} Behavioural
cloning~\citep{pomerleau1988alvinn} degrades under its own state
distribution, the failure DAgger~\citep{ross2011dagger} addresses. Our
COMPLETE($k$) family measures that degradation directly rather than inferring it
from a closed-loop gap: sweeping how many trailing actions the policy must
supply turns compounding error into a curve. We fine-tune with
PPO~\citep{schulman2017ppo} using GAE~\citep{schulman2016gae} and a KL anchor
back to the cloned policy.

\paragraph{Evaluation practice.} Our statistical treatment follows the
reproducibility literature in deep RL~\citep{henderson2018deeprl,
agarwal2021deeprl}, which documents how few-seed, few-episode comparisons
produce differences that do not survive resampling. We use paired
bootstrap intervals~\citep{efron1993bootstrap} because every policy faces
identical tasks. The point we add is complementary and, we think,
under-discussed: before asking whether two policies differ, one should ask what
the harness would score \emph{given a perfect policy}, because that number is
frequently not $1.0$.

\section{System}\label{sec:system}

\subsection{Action space and environment}

The environment wraps FreeCAD's~\citep{freecad} PartDesign workbench, exposed
through a Gym-style step interface~\citep{brockman2016gym}. A policy emits a triple
$(o, \mathbf{p}, t)$: an operation index over 38 operations, six continuous
parameters in $[0,1]^6$, and a target index selecting an edge, face or feature
from the live document. A codec decodes this into a validated structured action;
illegal actions are masked before sampling, and actions that pass the mask but
fail in the kernel are returned as failures rather than exceptions, so the
policy experiences them as negative outcomes rather than crashes.

\paragraph{Two interpreters.} FreeCAD ships its own Python, which cannot install
PyTorch. Rather than reimplement geometry, we split the system: the environment
runs under FreeCAD's interpreter, the policy under the system interpreter, and
they exchange newline-delimited JSON over a pipe. Everything in the environment
package is import-clean under both.

\subsection{Data}

An expert procedural generator produces designs across eight parametric families
(brackets, plates, flanges, spacers). Each design carries a natural-language
requirement, the full action trajectory, the resulting solid, and measured
properties. Requirements are generated \emph{from} the sampled parameters, so
ground truth is exact rather than annotated.

Crucially, expert trajectories are recorded as actions the codec can express.
This is not automatic; Section~\ref{sec:pitfalls} describes what happened when it
was not true. The current dataset round-trips \CodecSteps{} expert steps through
the codec with \CodecUnusable{} unusable and a worst-case drift of
\CodecWorstDrift{}~mm, two orders of magnitude below the $0.1$~mm tolerance every
constraint check uses.

\subsection{Policy and training}

The policy is a 1.14M-parameter model mapping requirement text and current
geometry to the next structured action, with separate heads for the operation,
the parameters and the target. We train by behavioural cloning for \BCEpochs{}
epochs, holding out by \emph{design} rather than by step --- a step-level split
puts actions from the same trajectory on both sides and inflates accuracy. We
then run PPO for \PPOIterations{} iterations with a clipped surrogate and a KL
anchor to the cloned policy, on a dense milestone reward.

\section{Benchmark}\label{sec:benchmark}

\subsection{Tasks}

KAIROS-CAD-v1 defines two task types over a frozen 70/15/15 split by design:

\begin{itemize}\itemsep2pt
\item \textbf{BUILD}: given only the requirement, build the part from an empty
      document.
\item \textbf{COMPLETE($k$)}: given the requirement and the expert's trajectory
      with its last $k$ actions removed, supply those $k$ actions. Larger $k$ is
      harder.
\end{itemize}

COMPLETE($k$) exists because BUILD alone is close to all-or-nothing at this scale
and reports $0.000$ for policies that differ substantially. Sweeping $k$
measures compounding error directly.

\subsection{Scoring}

Success rate alone cannot rank policies that all fail. We score a
\emph{prefix-ordered milestone ladder} with strictly dominating weights: opened a
sketch, drew geometry, made a solid, solid is valid, has a hole, all constraints
met, finished successfully. Credit stops at the first milestone missed. The
prefix rule is what stops a constraint check passing \emph{vacuously} on geometry
that was never built and ranking an empty document above a real but imperfect
part.

\subsection{Baselines}

Four baselines bound the task, and two of them audit the harness rather than
compete:

\begin{itemize}\itemsep2pt
\item \textbf{Oracle replay} pushes the ground-truth expert actions through the
      \emph{same codec} a policy uses. Its score is the ceiling.
\item \textbf{Scripted from spec} builds the one shape a fixed script can build
      from the parsed requirement: a plate with holes. It should nail flat
      families and fail bent ones.
\item \textbf{Immediate finish} terminates at once. Any policy below it is worse
      than doing nothing.
\item \textbf{Legal random} samples uniformly from the masked-legal actions.
\end{itemize}

\section{Results}\label{sec:results}

\subsection{Main results}

Table~\ref{tab:leaderboard} reports the leaderboard over the held-out test
split.

\begin{table}[t]
\centering
\begin{tabular}{lccccc}
\toprule
Policy & Progress & Success & Validity & Satisfaction & Efficiency \\
\midrule
Oracle replay\,$^\dagger$ & 0.594 & 0.500 & 0.924 & 0.656 & 1.000 \\
\midrule
Behavioural cloning & 0.445 & 0.281 & 0.957 & 0.453 & 0.620 \\
PPO (BC-anchored) & 0.413 & 0.250 & 0.951 & 0.477 & 0.675 \\
Scripted from spec & 0.224 & 0.000 & 1.000 & 0.211 & 0.640 \\
Immediate finish & 0.217 & 0.156 & 1.000 & 0.266 & 1.000 \\
Legal random & 0.146 & 0.000 & 0.676 & 0.268 & 0.604 \\
\bottomrule
\end{tabular}

\caption{KAIROS-CAD-v1, held-out test split.
$^\dagger$Oracle replay is not a policy: it replays the ground-truth expert
through the action codec, so its score is the \emph{ceiling}, and the gap
between it and $1.000$ measures what the codec and the target encoding cannot
express.}
\label{tab:leaderboard}
\end{table}

Two things are worth stating plainly. First, the oracle does not score
$1.000$; it scores \OracleCeiling{}. Reading behavioural cloning's
\BCProgress{} against $1.0$ overstates the shortfall by roughly a third of the
scale. Second, behavioural cloning reaches \BCHeldOutAccuracy{} next-action
accuracy on held-out designs while reaching \BCProgress{} in closed loop. The
gap is compounding error, and Table~\ref{tab:curve} measures it directly.

\begin{table}[t]
\centering
\begin{tabular}{lccccc}
\toprule
Policy & $k{=}1$ & $k{=}2$ & $k{=}4$ & $k{=}8$ & full build \\
\midrule
Oracle replay\,$^\dagger$ & 0.833 & 0.600 & 0.333 & 0.500 & 0.375 \\
Behavioural cloning & 0.833 & 0.600 & 0.333 & 0.000 & 0.000 \\
PPO (BC-anchored) & 0.833 & 0.600 & 0.000 & 0.000 & 0.000 \\
Immediate finish & 0.833 & 0.000 & 0.000 & 0.000 & 0.000 \\
\bottomrule
\end{tabular}

\caption{Success against difficulty. $k$ is the number of trailing actions the
policy must supply; the expert prefix is everything before it, so larger $k$ is
harder. Behavioural cloning holds up when one action remains and has collapsed
by eight.}
\label{tab:curve}
\end{table}

\subsection{The comparison that does not hold}

Behavioural cloning appears to edge PPO, \BCProgress{} against \PPOProgress{}.
Every policy faces identical tasks in identical order, which makes the paired
bootstrap the correct test: most of the variance at this sample size is
\emph{between tasks}, and pairing removes exactly that.

\begin{table}[t]
\centering
\begin{tabular}{lcccc}
\toprule
Comparison & $\Delta$ progress & 95\% CI & W/L/T & Separates \\
\midrule
\texttt{legal-random vs.\ oracle-replay} & -0.447 & $[-0.588, -0.310]$ & 2/27/3 & yes \\
\texttt{immediate-finish vs.\ oracle-replay} & -0.376 & $[-0.527, -0.235]$ & 1/25/6 & yes \\
\texttt{oracle-replay vs.\ scripted-spec} & +0.369 & $[+0.234, +0.507]$ & 22/9/1 & yes \\
\texttt{bc vs.\ legal-random} & +0.298 & $[+0.204, +0.402]$ & 27/2/3 & yes \\
\texttt{legal-random vs.\ ppo} & -0.267 & $[-0.371, -0.174]$ & 2/26/4 & yes \\
\texttt{bc vs.\ immediate-finish} & +0.228 & $[+0.141, +0.329]$ & 22/1/9 & yes \\
\texttt{bc vs.\ scripted-spec} & +0.220 & $[+0.126, +0.323]$ & 21/2/9 & yes \\
\texttt{immediate-finish vs.\ ppo} & -0.196 & $[-0.293, -0.115]$ & 1/22/9 & yes \\
\bottomrule
\end{tabular}

\caption{Paired bootstrap on the per-task difference in progress score, 5{,}000
resamples. An interval straddling zero means this benchmark cannot separate the
two policies at this sample size.}
\label{tab:comparisons}
\end{table}

The result (Table~\ref{tab:comparisons}) is that \textbf{BC and PPO do not
separate}: \BCvsPPODelta{} with a 95\% interval of $\BCvsPPOCI$, and
\BCvsPPOWLT{} in wins, losses and ties. Reporting ``BC beats PPO'' from the
point estimates would repeat an error we had already made once, when a
six-episode evaluation reported $0.500$ for a policy that scored $0.286$ over
fourteen.

\subsection{What the policy is actually using}

\begin{table}[t]
\centering
\begin{tabular}{lccc}
\toprule
Condition & Progress & $\Delta$ & Validity \\
\midrule
Behavioural cloning (intact) & 0.445 & --- & 0.957 \\
\quad requirement shuffled & 0.343 & -22.9\% & 0.955 \\
\quad requirement blanked & 0.374 & -15.9\% & 1.000 \\
\quad action mask removed & 0.339 & -23.9\% & 0.407 \\
\bottomrule
\end{tabular}

\caption{Ablations. Each row corrupts one input and re-runs the identical suite.}
\label{tab:ablations}
\end{table}

Table~\ref{tab:ablations} separates two claims that are easy to conflate.
Shuffling the requirement --- pairing a design with a different design's text ---
costs the policy substantially, so it is reading the requirement rather than
memorising a family prior. But removing the action mask collapses validity,
which means a large share of the policy's low invalid-action rate belongs to the
mask and not to the policy. Both facts follow from the same experiment and only
one of them is flattering.

\section{Four ways to measure the harness instead of the policy}\label{sec:pitfalls}

Each of the following produced a number we believed and reported. Each is now
covered by an automated audit.

\paragraph{1. An action space that cannot express its own demonstrations.}
Six of eight families sketched their outline as a single polygon carrying an
arbitrary vertex list. The codec can only express a polygon as a \emph{regular}
$n$-gon, so those steps were unrepresentable: cloning silently dropped them
(7.81\% of all expert steps, touching 829 of 1{,}080 designs) and an oracle
replaying the expert could not rebuild the part. The measured ``ceiling'' was an
artefact of the encoding, not a property of the task. Expanding each profile
into one line segment per edge drove unrepresentable steps to zero.

\paragraph{2. Normalisation that clips instead of failing.}
The codec maps each parameter into a fixed range. Values outside that range were
\emph{clipped}. A clipped value is not a rounding error; it is a different
action. A sketch offset of $89.9$~mm against a $\pm 50$~mm range decoded back as
$50$~mm, so the oracle built the feature $40$~mm from where the expert put it ---
with every step reporting success and zero invalid actions. This silently
corrupted 260 expert steps across 206 of 1{,}080 designs (19.1\%) while every
audit still called the trajectory fully representable. Encoding now
\emph{raises} on an out-of-range value, and ranges are named once and shared by
both directions; they had drifted apart, so an in-range pad length round-tripped
$61.9$~mm away having raised nothing at all.

\paragraph{3. A requirement the ground truth cannot satisfy.}
The constraint checker reads bounds from the requirement \emph{text}, not from
the float the generator used. Rounding a minimum to nearest states a bound the
part can miss: a $6.1866$~mm wall written as ``6.2~mm'' makes the expert violate
its own requirement while having built exactly the right part. A second cause
compounded it --- one family declared its minimum wall as the plate thickness
while its gusset rib, also a wall, could be thinner. Together these marked 139 of
1{,}080 designs (12.9\%) as failures. Minimums now round down, maximums round up,
and the stated bound is derived from the thinnest feature.

\paragraph{4. Model selection on the evaluation set.}
An early PPO run selected its best checkpoint on the same pool the evaluation
scored, and reported $0.286$ closed-loop success on requirements it had trained
on. The frozen three-way split in Section~\ref{sec:benchmark} exists because of
this. The corrected number was $0.000$ --- and the $0.000$ was itself misleading
for a different reason, being measured only on full builds, which is why
COMPLETE($k$) exists.

\paragraph{The common shape.} None of these raised an exception. Each produced a
plausible number, in the expected direction, of the expected magnitude. What
caught them was not testing that the code runs but \emph{measuring the harness
against a known answer}: replaying ground truth through the agent's own
machinery and asking why the score was not perfect. We would suggest that any
benchmark whose action space is a lossy encoding of its demonstrations should
report an oracle-replay ceiling as a matter of course.

\section{Limitations}\label{sec:limitations}

The learned results here are weak in absolute terms, and we state the ceiling
rather than the excuse. Full builds succeed rarely; the benchmark separates BC
and PPO from the null baselines but not from each other. Target selection ---
choosing which edge or face an operation applies to --- is not recoverable by the
current encoder, which caps families whose recipes pattern a specific feature.
The RL stage is single-seed. Families are procedural, so requirement language is
templated rather than natural. Wall thickness is measured by ray casting, which
gives an upper bound: failing it is conclusive, passing it is necessary but not
sufficient.

\section{Conclusion}

We built a CAD agent and a benchmark, and the most durable result was neither the
agent nor the score but the discipline of measuring the measuring apparatus. Four
separate defects in our pipeline each produced a confident wrong number, and in
every case the thing that exposed it was replaying a known-correct trajectory
through the agent's own machinery and refusing to accept a score below perfect
without an explanation. We release the environment, the dataset generator, the
frozen benchmark, and the audits.

\section*{Reproducibility}

Code, dataset generator, frozen splits, and every audit in this paper are at
\url{https://github.com/sunkalpchandra/kairos-cad}. Every number in every table
is generated from committed run artifacts by
\texttt{scripts/build\_paper\_tables.py}; none is typed by hand. The dataset is
regenerated deterministically from seeds by \texttt{make generate-data}.

\bibliographystyle{plainnat}
\bibliography{refs}

\end{document}
